\section{Propositonal Formulas}

We choose a semantic-central definition and define propositional formulas as maps from states to

\subsubsection{Formulas as maps}


We now choose a semantic centric approach to propositional logic, by defining formulas as boolean tensors.
Then we investigate the corresponding syntax of formulas as specification of a tensor network decomposition of the \basisEncoding{} of formulas.

% Intro of formulas
Logics is especially useful in interpreting boolean tensors representing Propositional Knowledge Bases, based on connections with abstract human thinking.
To make this more precise, we associate each such tensor with a formula $\exformula$ being a composition of the atomic variables with logical connectives as we proof next.

\begin{definition}
    \label{def:formulas}
    A propositional formula $\formulaat{\shortcatvariables}$ depending on $\atomorder$ atoms $\catvariableof{\atomenumerator}$ is a boolean-valued tensor
    \begin{align*}
        \formulaat{\shortcatvariables} \defcols \atomstates \rightarrow \ozset \subset \rr \, .
    \end{align*}
    We call a state $\shortcatindices \in \atomstates$ a model of a propositional formula $\formula$, if
    \begin{align*}
        \formulaat{\indexedshortcatvariables}=1 \, .
    \end{align*}
    If there is a model to a propositional formula, we say the formula is satisfiable.
\end{definition}

% Boolean Tensors
The propositional formulas coincide therefore with the boolean tensors (see \defref{def:booleanTensor}).


% Decomposition into model sums
Since propositional formulas are binary valued tensors, the generic decomposition of \lemref{lem:tensorBasisDecomposition} simplifies to
\begin{align}
    \label{eq:formulaModelDecomposition}
    \formulaat{\shortcatvariables} = \sum_{\catindices\in\atomstates} \formulaat{\indexedcatvariables} \cdot \onehotmapofat{\shortcatindices}{\shortcatvariables} \\
    = \sum_{\catindices\in\atomstates \wcols \formulaat{\indexedcatvariables}=1}  \onehotmapofat{\catindices}{\shortcatvariables} \, .
\end{align}
Thus, any propositional formula is the sum over the one-hot encodings of its models.
This is equal to the encoding of the set of models, which will be introduced in \charef{cha:basisCalculus} (see \defref{def:subsetEncoding}).

We depict this decomposition in the diagrammatic notation by
\begin{center}
    \input{./partI/0-propositional-logics/tikz/formula_direct.tex}
\end{center}

%% Semantic approach
We here chose a semantic approach to propositional logic in contrary to the standard syntactical approach.
Instead of defining formulas by connectives acting on atomic formulas, we define them here as binary valued functions of the states of a factored system.
They are interpreted by marking possible states as models, given the knowledge of $\exformula$.
The syntactical side will then be introduced later by studying decompositions of formulas.


\subsubsection{Knowledge bases}

Knowledge bases of propositional formulas are our first examples of tensor network decompositions.